\chapter{Detalle de infraestructura}\label{ap:infraestructura}

Este anexo amplía la sección~\ref{sec:implementacion-k8s} con los detalles de despliegue que conviene tener accesibles sin entrar al repositorio.

\section{Versión de OpenMetadata}\label{sec:version-openmetadata}

El despliegue de referencia del caso de uso de validación utiliza \textbf{OpenMetadata 1.12.9}, instalado mediante los \emph{charts} oficiales de Helm \cite{openmetadataKubernetes}. La versión exacta puede consultarse en cualquier momento contra la instancia viva mediante el endpoint \fqn{/api/v1/system/version} de la \ac{API} \ac{REST}. Las dependencias gestionadas por el \emph{chart} \fqn{openmetadata-dependencies} son MySQL como almacén de metadatos y OpenSearch como motor de búsqueda. Las entidades de la \ac{API} empleadas por la plataforma (servicios, bases de datos, esquemas, tablas, columnas, \emph{custom properties} y \emph{tags}) son estables en la rama 1.12.x, por lo que la lógica de gobierno no depende de detalles internos de una versión concreta.

\section{Matriz de versiones del sistema}\label{sec:matriz-versiones}

La reproducibilidad exige que todas las versiones del sistema estén fijadas y documentadas. La tabla~\ref{tab:matriz-versiones} recoge las versiones empleadas en el caso de uso de validación y el punto del repositorio donde quedan ancladas. El detalle operativo se mantiene actualizado en \fqn{docs/versiones.md}.

\begin{table}[h]
\footnotesize
\centering
\renewcommand{\arraystretch}{1.15}
\begin{tabular}{|L{4.1cm}|L{3.0cm}|L{4.7cm}|}
\hline
\rowcolor{tema!10}
\bft{Componente} & \bft{Versión} & \bft{Dónde se fija} \\ \hline
OpenMetadata (app y charts) & 1.12.9 & \texttt{-{}-version} en \fqn{launch_infra.ps1} \\ \hline
Helm (CLI) & 3.14.4 & \fqn{.tools/helm-v3.14.4/} \\ \hline
Kind & 0.31.0 & prerrequisito (\fqn{check_prereqs.ps1}) \\ \hline
kubectl & 1.34.1 & prerrequisito \\ \hline
Docker Engine & 28.5.1 & prerrequisito \\ \hline
Python & 3.12 (\textgreater{}= 3.10) & \fqn{tfm_ingestor/pyproject.toml} \\ \hline
Node.js & 22.11 & prerrequisito de \fqn{web/} \\ \hline
pnpm & 11.x & gestor Node canónico \\ \hline
Next.js / React / TypeScript & 16.2 / 19.2 / 6.0 & \fqn{web/package.json} \\ \hline
pyshacl & \textgreater{}= 0.28 & extra \texttt{validation} (\fqn{pyproject.toml}) \\ \hline
PyJWT / cryptography & \textgreater{}= 2.8 / 42 & extra \texttt{infra} (\fqn{pyproject.toml}) \\ \hline
fpdf2 & \textgreater{}= 2.7 & extra \texttt{report} (\fqn{pyproject.toml}) \\ \hline
Bundle SHACL DCAT-AP-ES & 1.0.0 (\texttt{f2c8a888}) & \fqn{resources/shacl/manifest.json} \\ \hline
Perfiles DCAT & DCAT-AP-ES 1.0.0\newline DCAT-AP 2.1.1\newline HVD 2.2.0 & documentación oficial \cite{datosGobDcatApEs2025} \\ \hline
\end{tabular}
\caption{Matriz de versiones fijadas del sistema. Elaboración propia.}
\label{tab:matriz-versiones}
\end{table}

\section{Valores Helm de OpenMetadata}

Los valores canónicos viven en \fqn{k8s/openmetadata.values.yaml} y \fqn{k8s/openmetadata-dependencies.values.yaml}. Documentan: tamaño de pods, configuración de Elasticsearch (OpenSearch), credenciales por defecto del JWT de servicio, namespace y puertos expuestos.

\section{Manifiesto de PostgreSQL de referencia}

El manifiesto \fqn{k8s/postgres-demo.yaml} crea un \texttt{Deployment}, un \texttt{Service} y un \texttt{ConfigMap} con el contenido de \fqn{sql/opendata_demo_init.sql}. La inicialización es automática: el contenedor monta el \ac{SQL} como volumen y lo ejecuta en el arranque.

\section{Diagrama de despliegue}

El diagrama de la figura~\ref{fig:kubernetes-reproducible} resume el despliegue. En operación real, el \texttt{kubectl port-forward} expone OpenMetadata en \fqn{localhost:8585} y PostgreSQL en \fqn{localhost:5432}.

\section{Notas operativas}

\begin{itemize}
    \item El clúster Kind se borra con \texttt{kind delete cluster --name tfm} si se requiere partir de cero.
    \item Las credenciales por defecto son las de la documentación oficial de OpenMetadata~\cite{openmetadataDocs}; en producción deberían rotarse.
    \item No se han activado \emph{políticas de red} ni \ac{RBAC} avanzado por estar fuera del alcance académico. Las \emph{políticas de red} (recurso \texttt{NetworkPolicy} de Kubernetes) restringen qué \emph{pods} pueden comunicarse entre sí y hacia el exterior; el \ac{RBAC} (\emph{Role-Based Access Control}) limita qué acciones puede ejecutar cada usuario o cuenta de servicio sobre los recursos del clúster. Ambos son controles de \emph{hardening} habituales en producción que no aportan valor al caso de uso de validación, centrado en gobierno de metadatos y reproducibilidad.
\end{itemize}

\section{Instalación portable y organismos con OpenMetadata propio}\label{sec:instalacion-portable}

La plataforma es portable y no exige desplegar su propio OpenMetadata. El despliegue con Kind, Helm y PostgreSQL de referencia descrito en este anexo solo construye un entorno autocontenido y reproducible para el caso de uso de validación. La lógica de gobierno reside íntegramente en el núcleo Python \texttt{om\_dcat\_sync}, que únicamente necesita acceso a la \ac{API} \ac{REST} de un OpenMetadata y un token \ac{JWT} válido.

En un organismo o empresa que ya disponga de OpenMetadata, la integración no requiere recrear infraestructura: basta con apuntar la plataforma a la instancia existente.

\begin{itemize}
    \item Configurar \fqn{OPENMETADATA_BASE_URL} hacia la \ac{API} del OpenMetadata del cliente (por ejemplo, \texttt{https://catalogo.organismo.es/api/v1}).
    \item Proporcionar un token de servicio en \fqn{OPENMETADATA_TOKEN} o \fqn{OPENMETADATA_JWT_TOKEN}, obtenido de un \emph{bot} de ingesta del propio OpenMetadata.
    \item Omitir los pasos de creación de clúster (\fqn{launch_infra.ps1}) y ejecutar directamente el workflow de gobierno, exportación y validación, ya sea desde el \ac{CLI} o desde la consola web.
\end{itemize}

De este modo, el mismo flujo (descubrimiento, curación funcional, sincronización repetible, exportación \ac{DCATAPES} y validación \ac{SHACL}) opera sobre el catálogo técnico que el cliente ya mantiene, sin imponer su motor de despliegue ni duplicar su gobierno. La compatibilidad se ha verificado contra OpenMetadata 1.12.x (sección~\ref{sec:version-openmetadata}). Las versiones próximas que conserven las mismas entidades de la \ac{API} \ac{REST} son igualmente compatibles. Esta separación entre lógica portable y despliegue de referencia es la que permite ofrecer la plataforma como capa de conformidad \ac{DCATAPES} sobre un OpenMetadata preexistente.
